\documentclass[12pt,letterpaper]{article}
\usepackage{graphicx,textcomp}
\usepackage{natbib}
\usepackage{setspace}
\usepackage{fullpage}
\usepackage{color}
\usepackage[reqno]{amsmath}
\usepackage{amsthm}
\usepackage{fancyvrb}
\usepackage{amssymb,enumerate}
\usepackage[all]{xy}
\usepackage{endnotes}
\usepackage{lscape}
\newtheorem{com}{Comment}
\usepackage{float}
\usepackage{hyperref}
\newtheorem{lem} {Lemma}
\newtheorem{prop}{Proposition}
\newtheorem{thm}{Theorem}
\newtheorem{defn}{Definition}
\newtheorem{cor}{Corollary}
\newtheorem{obs}{Observation}
\usepackage[compact]{titlesec}
\usepackage{dcolumn}
\usepackage{tikz}
\usetikzlibrary{arrows}
\usepackage{multirow}
\usepackage{xcolor}
\newcolumntype{.}{D{.}{.}{-1}}
\newcolumntype{d}[1]{D{.}{.}{#1}}
\definecolor{light-gray}{gray}{0.65}
\usepackage{url}
\usepackage{listings}
\usepackage{color}
\usepackage{etoolbox}
\makeatletter
\preto{\@verbatim}{\topsep=0pt \partopsep=0pt }
\makeatother


\definecolor{codegreen}{rgb}{0,0.6,0}
\definecolor{codegray}{rgb}{0.5,0.5,0.5}
\definecolor{codepurple}{rgb}{0.58,0,0.82}
\definecolor{backcolour}{rgb}{0.95,0.95,0.92}

\lstdefinestyle{mystyle}{
	backgroundcolor=\color{backcolour},   
	commentstyle=\color{codegreen},
	keywordstyle=\color{magenta},
	numberstyle=\tiny\color{codegray},
	stringstyle=\color{codepurple},
	basicstyle=\footnotesize,
	breakatwhitespace=false,         
	breaklines=true,                 
	captionpos=b,                    
	keepspaces=true,                 
	numbers=left,                    
	numbersep=5pt,                  
	showspaces=false,                
	showstringspaces=false,
	showtabs=false,                  
	tabsize=2
}
\lstset{style=mystyle}
\newcommand{\Sref}[1]{Section~\ref{#1}}
\newtheorem{hyp}{Hypothesis}

\title{Problem Set 1}
\date{Due: October 9, 2025}
\author{Ellen Beattie}

\begin{document}
	\maketitle
	
	\section*{Question 1: Education}

A school counselor was curious about the average of IQ of the students in her school and took a random sample of 25 students' IQ scores. The following is the data set:\\
\vspace{.25cm}

\lstinputlisting[language=R, firstline=41, lastline=41]{PS01.R}  

\vspace{0.5cm} 

\noindent
\subsection*{1.1: Finding a 90\% confidence interval for the average student IQ in school}
\vspace*{0.3cm}

\noindent
Step 1: Start by finding the point estimate (in this case average) from the school sample. 
\lstinputlisting[language=R, firstline=43, lastline=44]{PS01.R}\ 
	\begin{verbatim}
		[1] 98.44
	\end{verbatim} 
Step 2: Determine the variation in the sample (with standard deviation) and then the standard error.
(REASONS/EXPLANATION TBF)
\lstinputlisting[language=R, firstline=47, lastline=51]{PS01.R}\
 	\begin{verbatim}
 		[1] 13.09287
 		[1] 2.618575
 	\end{verbatim}
Step 3/4: Find relevant t-value given the confidence level and degrees of freedom in test
(REASONS/EXPLANATION TBF)
\lstinputlisting[language=R, firstline=55, lastline=57]{PS01.R}\
\begin{verbatim}
	[1] 1.317836
\end{verbatim}
Step 5: Use values to generate a confidence interval around the mean that includes a margin of error 
(REASONS/EXPLANATION TBF)
\lstinputlisting[language=R, firstline=60, lastline=63]{PS01.R}\
\begin{verbatim}
	[1]  94.98915 101.89085
\end{verbatim}
The confidence interval for the average student IQ in school is [94.99,101.89]. Meaning based on assumptions of the central limit theorem if 100 cases were sampled 90{\%} of the average scores would lie in this range. 
\noindent	
\vspace*{1cm}
	
\subsection*{1.2: Hypothesis test with $\alpha=0.05$ to determine whether the average student IQ in her school is higher than the average IQ score (100) among all the schools in the country}	
Step 1: Assumptions about data [TBF] . Sample size relatively small $n < 30$. \\
Step 2: Set up null and alternative hypothesis
Null hypothesis is the school (sample) average IQ score $\bar{y}$ is less than or equal to the national (population) average 100 IQ score. Alterative hypothesis is the sample average is greater than the population average 100. 
	\begin{equation}
		H_0: \bar{y} \le 100 ; H_a: \bar{y} > 100
	\end{equation} \\
	
Step 3: Calculate test statistic
Given small sample uses t-test to calculate statistic.
(more about THEORY TBF)
\lstinputlisting[language=R, firstline=69, lastline=71]{PS01.R}\
\begin{verbatim}
	[1] -0.5957439
\end{verbatim} 
(to be filled what test score means) \\

Step 4: Calculate p-value
(tbf p-value is)
\lstinputlisting[language=R, firstline=74, lastline=74]{PS01.R}\
\begin{verbatim}
	[1] 0.7215383
\end{verbatim} \

Step 5: Make conclusion
Since $p \ge 0.05$  
One-sided results
No evidence to reject null (tbc)

\newpage

	\section*{Question 2: Political Economy}

\noindent Researchers are curious about what affects the amount of money communities spend on addressing homelessness. The following variables constitute our data set about social welfare expenditures in the USA. \\
\vspace{.5cm}


\begin{tabular}{r|l}
	\texttt{State} &\emph{50 states in US} \\
	\texttt{Y} & \emph{per capita expenditure on shelters/housing assistance in state}\\
	\texttt{X1} &\emph{per capita personal income in state} \\
	\texttt{X2} &  \emph{Number of residents per 100,000 that are "financially insecure" in state}\\
	\texttt{X3} &  \emph{Number of people per thousand residing in urban areas in state} \\
	\texttt{Region} &  \emph{1=Northeast, 2= North Central, 3= South, 4=West} \\
\end{tabular}

\vspace{.5cm}
\noindent

\subsection*{2.0 Importing and Exploring the dataset \texttt{expenditure}.}
\vspace{.5cm}
\lstinputlisting[language=R, firstline=79, lastline=79]{PS01.R}  
\vspace{.5cm}


\subsection*{2.1 Relationships among \emph{Y}, \emph{X1}, \emph{X2}, and \emph{X3} }

\vspace{.5cm}

\subsection*{2.2 Relationship between \emph{Y} and \emph{Region}.}
\vspace{.5cm}

\subsection*{2.3 Relationship between \emph{Y} and \emph{X1} then including variable Region}


\end{document}
